% ===================================================================
%  TAVOLA N. 11 — La città e i suoi edifici. --- L'incendio.
%  Traduction 2026 d'apres texto/io, controlee sur texto/fr et
%  texto/en. Consigne : voir texto/it/10-tavola-01.tex.
%
%  LE TABLEAU 11 ALLEMAND AVAIT TROUVE QUATRE INSTITUTIONS PARLANT
%  CHACUNE LA LANGUE DU PAYS QUI L'AVAIT FONDEE. DEUX D'ENTRE ELLES
%  SONT ITALIENNES, ET ON LES VOIT ICI DE L'INTERIEUR.
%
%  LE SINDACO DE L'ALINEA 11 DE LA SECONDE SCENE (93). Le maire
%  s'appelle en italien sindaco, du grec « syndikos » —
%  « syn », avec, et « dike », le droit : CELUI QUI PLAIDE POUR LA
%  COMMUNAUTE. C'est un role de proces, non un rang.
%     francais  « maire »        < maior, le plus grand ;
%     anglais   « mayor »        le meme ;
%     allemand  « Bürgermeister » le maitre des bourgeois.
%  Les trois nomment un superieur ; l'italien nomme un AVOCAT. Et cela
%  se comprend : le premier officier du comune medieval etait celui
%  qui le representait en justice contre l'eveque ou contre
%  l'empereur. A cote de lui l'italien a garde le podestà,
%  « potestas », le juge etranger qu'on faisait venir d'une autre
%  ville pour qu'il n'ait aucun parent dans celle-ci. LES DEUX
%  MAGISTRATURES DE LA VILLE ITALIENNE SONT DEUX MAGISTRATURES
%  JUDICIAIRES.
%
%  LA BORSA DE L'ALINEA 5 (14). Elle n'est pas italienne du tout : le
%  mot vient de Bruges, de la famille Van der Beurse, dont l'hotel
%  portait trois BOURSES sur son ecu et devant lequel les marchands se
%  reunissaient. LE NOM DE LA BOURSE DE TOUTE L'EUROPE EST UN NOM DE
%  FAMILLE FLAMAND.
%
%  Or la banca de l'alinea 10 (32) est italienne, et le
%  tableau 11 allemand avait montre qu'elle etait le BANC du changeur.
%  Les deux institutions de l'argent sont donc sur la meme page, et
%  toutes deux portent le nom d'un MEUBLE : le banc ou l'on comptait,
%  la bourse ou l'on serrait. Ni l'une ni l'autre ne porte le nom de
%  ce qu'elle fait.
%
%  ET LES POMPIERS DE LA SECONDE SCENE SONT LE CONTRE-EXEMPLE. Le
%  tableau 11 allemand avait montre la Feuerwehr nommee d'un coup en
%  1846, avec tout son materiel. L'italien dit pompieri, mot
%  francais, de la POMPE — l'instrument, non le feu, non la defense.
%  Trois langues, trois choses vues dans le meme corps : l'allemand
%  voit la defense, le francais et l'italien voient la machine,
%  l'estonien voyait l'homme qui repousse. LA CHOSE EST NEUVE ET
%  CHACUN LA NOMME PAR LE COTE QU'IL REGARDE.
% ===================================================================

%%K t11-apar-1 apar
\VUsaut{2.0mm}
\VUcentre{13.2pt}{90}{TERZA SERIE}
\VUsaut{3.0mm}
\VUfilet{16mm}
\VUsaut{5.6mm}
\VUtitre{15.0pt}{60}{TAVOLA N. 11}
\VUsaut{1.4mm}
\VUpk{11.4pt}{40}{La città e i suoi edifici. --- L'incendio.}
\VUsaut{2.2mm}

%%K t11-tit-1 sub
\VUpk{10.2pt}{40}{I sobborghi.}
\VUsaut{3.0mm}

%%K t11-c1-01-1 p
1. --- Abito in una grande città a cento chilometri dal mare, che è però un
\VUgras{porto} \textsuperscript{(1)} di prim'ordine. Il fiume che le passa accanto è
qui largo quattrocento metri e forma una bellissima rada.

%%K t11-c1-02-1 p
2. --- Tutta la parte della città lungo le \VUgras{banchine} \textsuperscript{(2)} ha
un aspetto del tutto marittimo e industriale, e forma un \VUgras{sobborgo}
\textsuperscript{(3)} abitato solo da marinai e da operai. Questi lavorano nelle
\VUgras{fabbriche} \textsuperscript{(4)} e all'\VUgras{officina del gas}
\textsuperscript{(5)}, il cui alto \VUgras{camino} \textsuperscript{(6)} e il cui
\VUgras{gasometro} si vedono da lontano.

%%K t11-c1-03-1 p
3. --- Nel Medioevo la città era fortificata, e si trovano ancora dei resti delle sue
mura, in particolare la \VUgras{porta} \textsuperscript{(8)} del vecchio
\VUgras{castello} \textsuperscript{(7)} fiancheggiato dalle sue due \VUgras{torri}
\textsuperscript{(9)}, che oggi servono da \VUgras{carcere}.

%%K t11-c1-04-1 p
4. --- In tutto questo \VUgras{quartiere}, che ha un'aria molto antica, un solo
edificio è relativamente moderno: la \VUgras{dogana} \textsuperscript{(10)}. Sta fra la
banchina e la piccola \VUgras{via} \textsuperscript{(11)} che porta al vecchio
\VUgras{municipio} \textsuperscript{(12)}. Quest'ultimo edificio si riconosce facilmente
dal gigantesco \VUgras{campanile} \textsuperscript{(13)} che domina la città vecchia.

%%K t11-c1-05-1 p
5. --- Dall'altra parte della via \VUgras{sorge} un edificio nello stesso stile della
dogana: è la \VUgras{borsa} \textsuperscript{(14)}. Una delle sue facciate dà su una
piccola piazza dove si trova la \VUgras{prefettura} \textsuperscript{(15)}. La nostra
città non è una capitale, ma è comunque un importante capoluogo.

%%K t11-tit-2 sub
\VUsaut{5.0mm}
\VUpk{10.2pt}{40}{La parte alta della città.}
\VUsaut{3.4mm}

%%K t11-c1-06-1 p
6. --- Il presidio, piuttosto numeroso, è alloggiato in \VUgras{caserme}
\textsuperscript{(16)} ampie e molto salubri; arrivano fino alla cancellata dei
\VUgras{giardini pubblici} \textsuperscript{(17)}. Il \VUgras{viale}
\textsuperscript{(19)} che porta a questi giardini passa dietro la
\VUgras{cassa di risparmio} \textsuperscript{(18)} e sbocca sulla piazza del
\VUgras{duomo} \textsuperscript{(20)}.

%%K t11-c1-07-1 p
7. --- Tutti questi edifici si dispongono a gradini su una piccola collina dove sorge
anche il nostro grande \VUgras{liceo} \textsuperscript{(21)}.

%%K t11-c1-07-2 p
Scendendo per una strada in leggera pendenza che sbocca sul corso, si arriva al
\VUgras{palazzo di giustizia} \textsuperscript{(22)}, davanti al quale c'è un piacevole
\VUgras{giardinetto} \textsuperscript{(26)}. Questa \VUgras{piazza} non è molto grande,
ma in mezzo alla città fa un'oasi di verde e di frescura. Per questo i cittadini la
frequentano molto e vengono volentieri a sedersi sotto la tenda del \VUgras{caffè}
\textsuperscript{(25)} per ascoltare la banda militare che suona il giovedì e la
domenica nella \VUgras{cassa armonica} \textsuperscript{(27)} in mezzo ai prati e alle
aiuole.

%%K t11-c1-08-1 p
8. --- Su uno di quei prati sorge una \VUgras{statua} \textsuperscript{(28)} di bronzo,
un monumento eretto alla memoria di un nostro concittadino che ha reso grandi servizi
alla città in cui è nato.

%%K t11-tit-3 sub
\VUsaut{4.0mm}
\VUpk{10.2pt}{40}{Muoversi in città.}
\VUsaut{3.4mm}

%%K t11-c1-09-1 p
9. --- La nostra città non ha ancora la luce elettrica; ha solo i \VUgras{lampioni a
gas} \textsuperscript{(29)}. Ma la \VUgras{stampa locale} si batte perché
quell'illuminazione sia migliorata, \VUgras{così come} è già stata migliorata la
\VUgras{trazione} dei tram. Le vecchie \VUgras{vetture} tirate dai \VUgras{cavalli}
sono state sostituite da pratici \VUgras{tram elettrici} \textsuperscript{(31)}, che
portano i passeggeri rapidamente da un \VUgras{capo all'altro della città} a un prezzo
molto \VUgras{basso}.

%%K t11-c1-10-1 p
10. --- Accanto al palazzo di giustizia c'è la \VUgras{banca} \textsuperscript{(32)},
dove si scontano le cambiali commerciali. Gli \VUgras{esattori}
\textsuperscript{(33)} vanno a casa della gente a riscuotere quello che è dovuto.

%%K t11-tit-4 sub
\VUsaut{4.0mm}
\VUpk{10.2pt}{40}{Arte e scienza.}
\VUsaut{3.4mm}

%%K t11-c1-11-1 p
11. --- Il bel palazzo lì accanto è il \VUgras{museo}. Ospita insieme la
\VUgras{biblioteca} comunale \textsuperscript{(34)}, belle \VUgras{collezioni di
storia naturale} \textsuperscript{(35)} (un museo di storia naturale) e una graziosa
\VUgras{pinacoteca} \textsuperscript{(36)} che fa una splendida \VUgras{esposizione}
permanente.

%%K t11-c1-11-2 p
È aperto a tutti due volte alla settimana, ma i visitatori possono vederlo qualunque
giorno con una piccola mancia al \VUgras{custode} \textsuperscript{(37)}. Vengono a
lavorarci dei \VUgras{pittori} \textsuperscript{(38)}. Si vedono \VUgras{davanti} ai
loro cavalletti, con la \VUgras{tavolozza} in mano, che copiano i quadri famosi.

%%K t11-c1-12-1 p
12. --- Sulle alture dietro il museo si intravedono la \VUgras{cupola}
\textsuperscript{(40)} dell'\VUgras{ospedale} \textsuperscript{(39)} e gli edifici
della \VUgras{scuola normale} \textsuperscript{(41)}, dove si sono formati i maestri
delle nostre scuole elementari. In piazza del Teatro c'è un \VUgras{ufficio postale}
\textsuperscript{(43)} sistemato in una \VUgras{casa privata} \textsuperscript{(42)}.

%%K t11-tit-5 sub
\VUsaut{5.0mm}
\VUpk{11.4pt}{40}{L'incendio.}
\VUsaut{3.0mm}

%%K t11-tit-6 sub
\VUpk{10.2pt}{40}{Al fuoco!}
\VUsaut{3.4mm}

%%K t11-c2-01-1 p
1. --- Una notizia allarmante corre per la città: è scoppiato un incendio al teatro!
Quando arrivo alla \VUgras{piazza} \textsuperscript{(96)}, la folla si è già radunata
sul \VUgras{marciapiede} \textsuperscript{(46)} e sulla \VUgras{carreggiata}
\textsuperscript{(47)}. Gli \VUgras{impiegati delle poste} \textsuperscript{(44)} e i
\VUgras{portalettere} \textsuperscript{(45)} escono dall'ufficio postale. Un
\VUgras{tranviere} \textsuperscript{(48)} ha dovuto fermare il suo tram per far passare
un'\VUgras{ambulanza} \textsuperscript{(50)} municipale che arriva con un
\VUgras{chirurgo} \textsuperscript{(52)} e un'\VUgras{infermiera}
\textsuperscript{(51)} per soccorrere un \VUgras{ferito} \textsuperscript{(53)}
portato su una \VUgras{barella} \textsuperscript{(54)} in un punto vicino
all'\VUgras{edicola dei giornali} \textsuperscript{(55)}.

%%K t11-c2-02-1 p
2. --- Una compagnia di fanteria di ritorno dall'esercitazione si è fermata per aiutare
a tenere l'ordine insieme alle \VUgras{guardie municipali} \textsuperscript{(61)} a
cavallo. I \VUgras{cavalieri}, i cui cavalli si impennano, ci riescono meglio dei
\VUgras{soldati} \textsuperscript{(57)} o dei \VUgras{carabinieri}
\textsuperscript{(63)}: spingono indietro i \VUgras{curiosi} \textsuperscript{(56)}. Del
resto i carabinieri hanno molto da fare. Uno di loro dice agli \VUgras{agenti}
\textsuperscript{(64)} dove portare un altro \VUgras{ferito} \textsuperscript{(65)}, un
coraggioso pompiere caduto da una scala.

%%K t11-c2-03-1 p
3. --- L'\VUgras{ufficiale} \textsuperscript{(66)} indica esattamente dove va messa la
\VUgras{pompa a vapore} \textsuperscript{(67)} che sta arrivando. In attesa di quella
macchina potente, ha fatto montare una \VUgras{pompa a mano} \textsuperscript{(68)}, il
cui lungo \VUgras{tubo} \textsuperscript{(69)} versa torrenti d'acqua sul fuoco. Sei
pompieri la manovrano, mentre il loro compagno, che tiene la \VUgras{lancia}
\textsuperscript{(70)}, dirige il \VUgras{getto} \textsuperscript{(71)} nel cuore
dell'incendio.

%%K t11-c2-04-1 p
4. --- Dall'altra parte del teatro è stata alzata la grande \VUgras{scala}
\textsuperscript{(72)}. Permette ai pompieri di avvicinarsi alle fiamme che salgono in
mezzo a volute di \VUgras{fumo} nero \textsuperscript{(75)}.

%%K t11-tit-7 sub
\VUsaut{5.0mm}
\VUpk{10.2pt}{40}{Atti di coraggio.}
\VUsaut{3.4mm}

%%K t11-c2-05-1 p
5. --- L'\VUgras{incendio} \textsuperscript{(74)} è scoppiato nel ridotto del
\VUgras{teatro} \textsuperscript{(73)}, durante una prova dell'\VUgras{opera} la cui
prima \VUgras{rappresentazione} doveva essere domani. Per fortuna in sala c'erano solo
pochi \VUgras{spettatori} \textsuperscript{(76)}. I poveretti si sono rifugiati sul
\VUgras{balcone} \textsuperscript{(77)}, dove coraggiosi \VUgras{pompieri}
\textsuperscript{(86)} vengono in loro aiuto con una scala e delle corde.

%%K t11-c2-06-1 p
6. --- Sotto il \VUgras{portico} \textsuperscript{(78)}, dove il \VUgras{cartellone}
\textsuperscript{(79)} annuncia lo spettacolo, si vedono fuggire i \VUgras{musicisti}
\textsuperscript{(81)} dell'\VUgras{orchestra} \textsuperscript{(80)}, che portano via
i loro strumenti: \VUgras{violino} \textsuperscript{(81)}, \VUgras{flauto}
\textsuperscript{(82)}, \VUgras{violoncello} \textsuperscript{(83)},
\VUgras{trombone} \textsuperscript{(84)}, \VUgras{tamburo} \textsuperscript{(85)} e
così via. Si cerca di salvare una parte dell'\VUgras{arredamento}
\textsuperscript{(87)}.

%%K t11-c2-07-1 p
7. --- Gli \VUgras{attori} \textsuperscript{(89)} e le \VUgras{attrici}
\textsuperscript{(88)}, \VUgras{cantanti} e cantatrici, sono dovuti scappare fuori nei
loro \VUgras{costumi di scena} \textsuperscript{(90)}. Il tenore spiega a un
\VUgras{giornalista} \textsuperscript{(92)} com'è successo. \VUgras{Il cronista} è
accorso fin dal primo momento insieme alle autorità: il \VUgras{prefetto}
\textsuperscript{(91)}, il \VUgras{sindaco} \textsuperscript{(93)}, il
\VUgras{questore} \textsuperscript{(94)} e un \VUgras{magistrato}
\textsuperscript{(95)}, che aprirà un'inchiesta per stabilire di chi è la
responsabilità.
\VUsaut{6.0mm}
\VUfilet{22mm}
